We check the Fourier normalization of the operator on this line.
Use normalized induction, so that
\[
 v\left(\begin{pmatrix}a&x\\0&d\end{pmatrix}g\right)
 =\mu_1(a)\mu_2(d)|a/d|_p^{1/2}v(g).
\]
Put $K_0(p)=\{g\in\mathrm{GL}_2(\Z_p):g_{21}\in p\Z_p\}$.
Use its character $k\mapsto\mu_2(k_{22})$.
This is the lower-right-character convention. Casselman's
upper-left-character convention is carried to it by
$h_p=\left(\begin{smallmatrix}0&1\\-p&0\end{smallmatrix}\right)$,
which normalizes $K_0(p)$ and interchanges its diagonal entries
under conjugation. We label the inducing characters so that
$\mu_1$ is unramified, using the interchange isomorphism for
irreducible principal series.
Normalize Haar measure by $\operatorname{vol}(K_0(p))=1$.
The character-weighted double-coset operator for
$K_0(p)\operatorname{diag}(p,1)K_0(p)$, with value one at
$\operatorname{diag}(p,1)$, is
\[
 H_p=\sum_{b\in\Z/p\Z}
       \Pi_{f,p}\left(\begin{pmatrix}p&b\\0&1\end{pmatrix}\right).
\]
Take the newvector supported on $B K_0(p)$ with $v(1)=1$.
Evaluating at $1$ gives
\[
 (H_pv)(1)=p\mu_1(p)p^{-1/2}=\sqrt p\,\mu_1(p).
\]
This is the Fourier operator $U_p^{(2)}$. Indeed, in the
Whittaker model use an additive character of conductor $\Z_p$.
Writing $a(n)=\operatorname{diag}(n,1)$ and
$n(x)=\left(\begin{smallmatrix}1&x\\0&1\end{smallmatrix}\right)$,
\[
 a(n)\begin{pmatrix}p&b\\0&1\end{pmatrix}=n(nb)a(np)
 \quad(n\in\Z_p),\qquad
 (H_pW)(a(n))=pW(a(np)).
\]
Normalize the archimedean Whittaker function by
$W_\infty(a(y))=y\exp(-2\pi y)$ for $y>0$.
Comparison of its positive-$n$ term with $y f(iy)$ gives the
finite Whittaker coefficient $a_n/n$. The displayed equality
therefore sends $a_n$ to $a_{pn}$. Hence
$U_p^{(2)}=\sqrt p\,\mu_1(p)$ on the new line, exactly the
arithmetic-Frobenius eigenvalue on the unramified Galois line.
This scalar is algebraic; the square roots only express the
unitary normalization. Old multiplicities at primes other than
$p$ carry the same scalar. Casselman's oldspace dimension formula
\cite[Corollary, p.~306, and Theorem~4, p.~313]{casselman}
does not introduce an additional multiplicity at this conductor-one
prime. For Eisenstein constituents, the two Galois characters and
the divisor-sum Fourier coefficients give the same equality directly;
Carayol's cuspidal theorem is not applied to them.

Choose an inertia element $\tau$ on which the quadratic character
is $-1$, and a Frobenius lift $\phi$. On a ramified block, a basis
of the two character lines gives
\[
 \rho(\tau)=\begin{pmatrix}1&0\\0&-1\end{pmatrix},\qquad
 \rho(\phi)=\begin{pmatrix}\lambda&0\\0&\mu\end{pmatrix},
 \qquad U_p^{(2)}=\lambda,\quad d(\phi)=\lambda\mu.
\]
For $\rho(g)=\left(\begin{smallmatrix}a&b\\c&e\end{smallmatrix}\right)$,
the two inertia differences are
\[
 t(g\tau)-t(g)=-2e,\qquad
 t(g\tau\phi)-t(g\phi)=-2e\mu.
\]
It follows in characteristic zero that
\[
 U_p^{(2)}\{t(g\tau\phi)-t(g\phi)\}
       =d(\phi)\{t(g\tau)-t(g)\}.                       \tag{O.12}
\]
On a block coming from level prime to $p$, the representation is
unramified and both differences vanish. This proves the identity
on its entire old multiplicity space even if $U_p^{(2)}$ is not
scalar or semisimple on that space. It therefore proves (O.12)
on the actual sum containing $F_U$.

These are identities on characteristic-zero forms before any
reduction. To relate them to the integral trace module above, choose
finitely many Fourier coordinates that inject its ambient space.
Every generator of that module has integral coordinates; hence the
module is a submodule of a finite free module over the complete
coefficient DVR. It is finite and closed. Simultaneous Chebotarev
approximation in the finitely many Galois representations makes
$t(g)$ and $d(g)$ limits of good Hecke and diamond operators on this
space. They preserve this integral module, so the identity can be
evaluated on the integral $F_U$ and then reduced. This does not
require integral projectors onto constituents. In particular no
division by the two in the displayed inertia differences is used.


